\documentclass[11pt]{ltjsarticle}
\usepackage{luatexja-fontspec}
\usepackage{amsmath,amssymb}
\usepackage[margin=22mm]{geometry}
\usepackage{graphicx}
\usepackage{longtable,booktabs,array}
\usepackage{calc}
\usepackage{xcolor}
\usepackage{hyperref}
\hypersetup{colorlinks=true,linkcolor=blue,urlcolor=blue,citecolor=blue}
\makeatletter
\def\maxwidth{\ifdim\Gin@nat@width>\linewidth\linewidth\else\Gin@nat@width\fi}
\makeatother
\setkeys{Gin}{width=\maxwidth,keepaspectratio}
\setcounter{secnumdepth}{0}
\setlength{\parindent}{0pt}
\setlength{\parskip}{0.5em}
\providecommand{\tightlist}{\setlength{\itemsep}{0pt}\setlength{\parskip}{0pt}}
\providecommand{\pandocbounded}[1]{#1}
\providecommand{\real}[1]{#1}
\newcounter{none}
\begin{document}
\section{Appendix: The Complex Structure is a Consequence, Not a
Posit}\label{appendix-the-complex-structure-is-a-consequence-not-a-posit}

\subsection{--- How the Complex Plane Arises from Two Real
Components}\label{how-the-complex-plane-arises-from-two-real-components}

\textbf{Author:} Noriaki Kihara (ORCID: 0009-0004-6753-4020)
\textbf{Date:} 18 June 2026 \textbf{DOI (this version):}
10.5281/zenodo.20742985 \textbf{Concept DOI (all versions):}
10.5281/zenodo.20742984 \textbf{Zenodo:}
https://zenodo.org/records/20742985 \textbf{Related paper:} ``On the
Logarithmic Map of Scale and the Separation of Scale and Quantization in
a High-Symmetry Limit,'' §4, §7.4 (Concept DOI:
\href{https://doi.org/10.5281/zenodo.20740841}{10.5281/zenodo.20740841})

\begin{center}\rule{0.5\linewidth}{0.5pt}\end{center}

\subsection{Purpose}\label{purpose}

This appendix develops independently a point compressed into a few lines
in §4 and the note in §7.4 of the related paper --- namely that,
\textbf{without positing the imaginary unit i in the axioms, the complex
structure arises as a consequence of two real components and one
rotation}. It adds no new claim to the parent paper; its aim is to make
explicit ``where i comes from,'' which is easily misread.

Ordinarily, quantum theory gives the complex Hilbert space as the basic
framework and posits the imaginary unit i at the outset. In this system
i is not placed in the axioms. What is placed in the axioms is only two
independent real conserved quantities --- the sum of squares (the
binding of difference, g) and the product (the edge of being/non-being,
ω) --- and the rotation connecting them. The complex structure (the
complex plane, i, the complex inner product) arises as a consequence of
this real two-ness.

This appendix is not a rigorous proof but an observation making explicit
the structure of the parent paper. No physical identification is made.
It does not claim to ``have explained why quantum theory is described by
complex numbers''; it only shows a path by which the complex structure
formally appears from the minimal axioms of this system.

\begin{center}\rule{0.5\linewidth}{0.5pt}\end{center}

\subsection{1. Starting point: two real components, no
i}\label{starting-point-two-real-components-no-i}

The axioms of this system close over the reals. Complex numbers are not
introduced. The basic objects are two independent real conserved
quantities (Axioms 4 and 5 of the parent paper).

Write the two real components as ν\textsubscript{1}, ν\textsubscript{2}. These are anonymous, and which
is called first or second is arbitrary (the anonymity of the parent
paper). The two form a pair of equal real axes.

The imaginary unit i does not exist at this stage. What is shown below
is that on top of the two real components ν\textsubscript{1}, ν\textsubscript{2}, two conserved
quantities --- the sum of squares and the area --- stand, and that they
form the same algebraic structure as the real and imaginary parts of a
complex number.

\begin{center}\rule{0.5\linewidth}{0.5pt}\end{center}

\subsection{2. The two conserved quantities: the sum of squares (g) and
the area
(ω)}\label{the-two-conserved-quantities-the-sum-of-squares-g-and-the-area-ux3c9}

For the two real components ν\textsubscript{1}, ν\textsubscript{2}, consider two independent conserved
quantities.

\textbf{Sum of squares (the binding of difference, g):}

\[g = \nu_1^2 + \nu_2^2.\]

This binds the two on equal footing and is invariant under relabeling
(ν\textsubscript{1}↔ν\textsubscript{2}). It is the invariant of SO(2) --- the rotation preserving the
sum of squares (parent paper, Axiom 5).

\textbf{Area (the edge of being/non-being, ω):} the product of
fluctuations

\[\omega = \delta_1\delta_2 = \tfrac{1}{2} \quad(\text{area / symplectic form; parent paper, Axiom 4}).\]

This is an area (a quantity spanned by the two axes), related to SO(1,1)
--- the hyperbolic, product-preserving transformation --- and an
invariant preserved under rotation (parent paper, Axiom 4). Note that
the product of the components themselves, ν\textsubscript{1}ν\textsubscript{2}, is a different bilinear
quantity that depends on the rotation θ, and is not ω (δ\textsubscript{1}δ\textsubscript{2}) itself
(§3).

By the note of the parent paper (mutual underivability of Axioms 4 and
5), g and ω cannot be derived from each other. ``Whether it is''
(product, ω) and ``whether it differs'' (sum of squares, g) are two root
distinctions standing independently under anonymity.

\begin{center}\rule{0.5\linewidth}{0.5pt}\end{center}

\subsection{3. The rotation: introducing cos and
sin}\label{the-rotation-introducing-cos-and-sin}

Represent the two real components by a radius ν and an angle θ. This is
merely writing the two components of a plane in polar form, not a new
assumption:

\[\nu_1 = \nu\cos\theta, \qquad \nu_2 = \nu\sin\theta.\]

cos and sin are parameters of the rotation (SO(2)) preserving the sum of
squares. Here, by cos²θ+sin²θ=1,

\[g = \nu_1^2 + \nu_2^2 = \nu^2(\cos^2\theta + \sin^2\theta) = \nu^2,\]

so the sum of squares g equals the square of the radius ν², and is
invariant under the rotation θ (determined by the radius alone). On the
other hand, the product of the individual components is

\[\nu_1\nu_2 = \nu^2\cos\theta\sin\theta = \tfrac{1}{2}\nu^2\sin 2\theta,\]

which varies with θ. This is a component-level bilinear quantity and is
not ω itself. The area ω (δ\textsubscript{1}δ\textsubscript{2}=1/2, §2) corresponding to the imaginary
part of the complex structure is a symplectic invariant preserved under
rotation.

\begin{center}\rule{0.5\linewidth}{0.5pt}\end{center}

\subsection{4. The arising of the complex
structure}\label{the-arising-of-the-complex-structure}

Now bind the two real components (ν\textsubscript{1}, ν\textsubscript{2}) into a single object. Write
the bound object as z, and call its two components the ``real part'' and
``imaginary part'':

\[z = \nu_1 + i\,\nu_2 = \nu(\cos\theta + i\sin\theta) = \nu\,e^{i\theta}.\]

The i that appears here is not posited. \textbf{i is the sign of the
90-degree rotation J connecting the two real axes} (the cos↔sin of §3).
The two-ness (two independent axes) gives the ``plane'' on which i acts,
and the 90-degree rotation J on that plane corresponds to i. That
\textbf{J²=−1} --- a 90-degree rotation done twice is a 180-degree
rotation, i.e.~multiplication by −1 --- becomes the rule i²=−1. i is not
a new number introduced from outside but a symbol writing the rotation
on the two axes algebraically. (Two independent axes alone are merely
the real plane $\mathbb{R}$²; only by adding the rotation J does it become the
complex plane $\mathbb{C}$.)

Under this binding, three correspondences stand:

\begin{itemize}
\tightlist
\item
  \textbf{Sum of squares g = ν\textsubscript{1}²+ν\textsubscript{2}² = \textbar z\textbar²} --- the
  square of the absolute value of the complex number (real part).
\item
  \textbf{Area ω = δ\textsubscript{1}δ\textsubscript{2}} --- a quantity belonging to the imaginary side
  of the complex number (symplectic form).
\item
  \textbf{Rotation cos↔sin} --- the complex structure J (90-degree
  rotation, corresponding to multiplication by i).
\end{itemize}

That is, the structure written entirely in reals --- two real
components, the sum of squares, the area, the rotation --- corresponds
one-to-one with the complex number z=νe\^{}\{iθ\} and its algebra
(absolute value, argument, multiplication by i). The complex plane is an
image that arises on top of the real two-ness (Figure 1).

\begin{figure}
\centering
\pandocbounded{\includesvg[keepaspectratio]{fig3_complex_from_two_real.pdf}}
\caption{Figure 1: The complex structure arising from two real
components. The point z rides on the real axes ν\textsubscript{1}, ν\textsubscript{2}; the radius ν=√g
is invariant under rotation (dashed circle), while the individual
components vary with the angle θ. The pair (ν\textsubscript{1}, ν\textsubscript{2}) corresponds to the
complex number z=νe\^{}\{iθ\}, and i appears as the sign of the
90-degree rotation J (J²=−1) connecting the two axes.}
\end{figure}

\textbf{Correspondence with the complex inner product.} Binding the two
conserved quantities into a single complex quantity gives the form

\[\langle\,\cdot\,,\cdot\,\rangle = g + i\,\omega,\]

where the real part g (sum of squares, distance, metric) and the
imaginary part ω (area, symplectic form) correspond to the real and
imaginary parts of the complex inner product. The g/ω separation of §7
of the parent paper states that the real and imaginary parts of this
complex inner product follow separate fates under the logarithmic map (g
is flattened, ω is invariant).

\begin{center}\rule{0.5\linewidth}{0.5pt}\end{center}

\subsection{5. Summary of the structure}\label{summary-of-the-structure}

{\def\LTcaptype{none} % do not increment counter
\begin{longtable}[]{@{}
  >{\raggedright\arraybackslash}p{(\linewidth - 2\tabcolsep) * \real{0.5000}}
  >{\raggedright\arraybackslash}p{(\linewidth - 2\tabcolsep) * \real{0.5000}}@{}}
\toprule\noalign{}
\begin{minipage}[b]{\linewidth}\raggedright
Structure written in reals
\end{minipage} & \begin{minipage}[b]{\linewidth}\raggedright
Correspondence in complex representation
\end{minipage} \\
\midrule\noalign{}
\endhead
\bottomrule\noalign{}
\endlastfoot
Two real components ν\textsubscript{1}, ν\textsubscript{2} & real and imaginary parts of the complex
number z \\
The 90-degree rotation J connecting the two axes & i (J²=−1,
i.e.~i²=−1) \\
Sum of squares g = ν\textsubscript{1}²+ν\textsubscript{2}² & \\
Area ω = δ\textsubscript{1}δ\textsubscript{2} & imaginary side (symplectic form) \\
Rotation cos↔sin (SO(2)) & the complex structure J (multiplication by i,
argument θ) \\
Binding of the two conserved quantities & complex inner product $\langle$·,·$\rangle$ =
g + iω \\
\end{longtable}
}

Complex numbers are not placed in the axioms. What is placed is only the
two independent real conserved quantities (the sum of squares g, the
area ω) and the rotation. i appears, at the stage of binding, as the
symbol of the 90-degree rotation. The complex plane / complex inner
product is a consequence of this real two-ness (the two-ness given by
the parent paper's ``limit of anonymity = 2'') and the rotation.

Why base 2, why two --- that is by Axiom 2 of the parent paper (the
limit of anonymity = 2). That the minimum and maximum at which
difference stands in the interior is two is the ground for ``two real
components,'' and also the ground for the complex structure being
two-dimensional (one real part, one imaginary part). The complex number
being two-dimensional is a reflection of the limit of anonymity being 2.

\begin{center}\rule{0.5\linewidth}{0.5pt}\end{center}

\subsection{6. Limitations (reading
notes)}\label{limitations-reading-notes}

This appendix makes explicit the structure of §4 and §7.4 of the parent
paper and adds no new claim.

\begin{enumerate}
\def\labelenumi{\arabic{enumi}.}
\item
  \textbf{This appendix does not claim to have explained ``why quantum
  theory is described by complex numbers.''} It only shows one path by
  which the complex structure formally appears from the minimal axioms.
  Isomorphism with the complexity of physical theory is a separate task.
\item
  \textbf{i appears as the symbol of the 90-degree rotation J (J²=−1).}
  This is not a claim that ``i has been given physical reality,'' but a
  consequence of the algebraic notation that binds the rotation on two
  real components into a single complex quantity. Two independent axes
  alone are merely the real plane $\mathbb{R}$²; only by adding the rotation J does
  it become the complex plane $\mathbb{C}$. Whether one writes the real two-ness
  and the rotation in complex notation or in two real components plus a
  rotation is a choice of representation; the content of this system
  closes over the reals.
\item
  \textbf{The mutual underivability of the sum of squares g and the area
  ω is a premise.} As in the note of the parent paper, g and ω are two
  independent root distinctions, and one is not derived from the other.
  The complex structure appears on top of these two standing
  independently.
\item
  \textbf{No physical identification is made.} The physical names of the
  two real components ν\textsubscript{1}, ν\textsubscript{2} are anonymous labels, and this appendix
  makes no identification with a particular physical quantity.
\item
  \textbf{The ground for two-ness depends on Axiom 2 of the parent
  paper.} ``Why two real components (why the complex number is
  two-dimensional)'' is attributed to the limit of anonymity = 2 (Axiom
  2). The derivation of that axiom itself is outside the scope of the
  parent paper and this appendix.
\end{enumerate}

\begin{center}\rule{0.5\linewidth}{0.5pt}\end{center}

\subsection{Reference}\label{reference}

\begin{itemize}
\tightlist
\item
  Kihara, N. ``On the Logarithmic Map of Scale and the Separation of
  Scale and Quantization in a High-Symmetry Limit --- A Limited
  Verification Report from a Minimal Axiom System,'' §4, §7.4 (the
  complex structure as a consequence of real two-ness), Axioms 4 and 5.
  Concept DOI:
  \href{https://doi.org/10.5281/zenodo.20740841}{10.5281/zenodo.20740841}.
\end{itemize}
\end{document}
